
% =============================================================
%  Chapter 2 — Background
% =============================================================
%    §2.1  LLM code generation and the BaxBench benchmark
%    §2.2  Containerised deployment and the limits of single-host evaluation
%    §2.3  Kubernetes and container orchestration
%    §2.4  Deployment configuration as a performance lever
%    §2.5  Functional correctness versus runtime performance
%    §2.6  Measuring performance under load
%    §2.7  Iterative refinement as an agentic improvement pattern
%    §2.8  Conceptual model: code and deployment as coupled levers
% =============================================================

\chapter{Background}
\label{ch:background}

Chapter~\ref{ch:introduction} motivated the problem this thesis addresses and
stated its central research question. This chapter introduces the technical
concepts and assumptions needed to follow the remainder of the thesis. It is
selective rather than encyclopaedic: each concept is included because it is
used to define the benchmark, interpret a result, or understand a design
choice. Chapter~\ref{ch:related-work} instead compares this work with prior
systems.

Section~\ref{sec:background-codegen} introduces \textsc{BaxBench}, the
code-generation benchmark this thesis extends, and briefly explains the role
of \textsc{AutoBaxBuilder} in constructing scenarios. Section
\ref{sec:background-containerization} explains why an isolated container is
not a model of a deployed service. Sections~\ref{sec:background-kubernetes}
and~\ref{sec:background-deployment-levers} introduce Kubernetes and the
deployment decisions exposed to the agent. Sections
\ref{sec:background-correctness-vs-performance} and
\ref{sec:background-goodput} distinguish correctness from runtime performance
and define the load-testing vocabulary used later. Section
\ref{sec:background-iterative-refinement} describes the general
feedback-driven improvement loop, and Section~\ref{sec:background-synthesis}
states the conceptual model that connects code, deployment, workload, and
cluster capacity.


% ----- §2.1  BaxBench -----
\section{BaxBench and LLM Code Generation}
\label{sec:background-codegen}

Large language models can now produce complete, multi-file backend
applications rather than isolated functions or patches. Evaluating such
systems requires an application-level contract: a description of the
behaviour to implement and executable checks that determine whether the
resulting service satisfies it.

\textsc{BaxBench}~\cite{vero2025baxbench} is the benchmark this thesis builds
on. It defines 392 tasks across 28 \emph{scenarios}. A scenario describes a
fixed backend workload through an API contract; a recipe-sharing service with
upload, comment, and rating endpoints is one example. Each scenario is paired
with 14 \emph{frameworks}, combinations of a programming language and a web
or application framework, spanning six languages. The same scenario can
therefore be implemented and compared across different technology stacks.

For a scenario--framework pair, a model generates a complete application.
\textsc{BaxBench} then applies two separate forms of checking. A
\emph{functional test suite} exercises the API contract end to end and checks
that responses have the expected behaviour. A separate \emph{security test
suite} probes known vulnerability classes, such as SQL injection or broken
authentication. The distinction matters for this thesis: functional tests
provide the correctness gate used before performance measurement, whereas
security testing is part of the original benchmark but is not reconstructed
as a second Kubernetes evaluation in this work.

\textsc{AutoBaxBuilder}~\cite{vonarx2025autobaxbuilder} addresses the cost of
constructing such benchmark inputs. Its LLM-driven construction pipeline
produces scenario descriptions, API contracts, functional tests, and
calibrated reference implementations; its full form can also construct
security checks. This thesis builds its offline scenarios on that pipeline and
extends it with load-test generation and verification. The inherited stages,
the adaptations made here, and the new workload stage are described in
Section~\ref{sec:methods-scenario-gen}; the details are not needed to
understand the benchmark's deployment model.

The original \textsc{BaxBench} evaluation builds each candidate into a
container and executes it alone on a single host. That environment is
appropriate for asking whether generated code satisfies an application
contract. The next section explains what it cannot reveal about the same code
once it becomes a replicated service.


% ----- §2.2  Containerization limits -----
\section{Containerised Deployment and the Limits of Single-Host Evaluation}
\label{sec:background-containerization}

A container packages an application with its user-space runtime dependencies,
such as its interpreter, libraries, and configuration, into a portable unit.
This makes it possible to build and execute many scenario--framework
combinations without provisioning a bespoke environment for each one. A
container improves packaging consistency, but it does not make performance
independent of the host: CPU, memory, storage, networking, and contention
still affect the running process.

Running one container on one host answers a narrow question: does one instance
of the application behave correctly for the tested request paths? It does not
answer questions that arise when the application is deployed as a service:
whether enough replicas exist for concurrent demand, whether replicas compete
for shared worker capacity, whether placement creates a hotspot, whether a
shared database becomes the bottleneck, or whether a connection pool sized
for one instance remains appropriate when several instances issue requests at
once. These are properties of the application \emph{as deployed}, not of its
source code in isolation.

The deployment setting considered in this thesis therefore adds at least
three elements absent from the isolated benchmark: multiple application
replicas, shared cluster resources, and a load generator that creates
overlapping requests. Kubernetes supplies the orchestration layer for that
setting.


% ----- §2.3  Kubernetes -----
\section{Kubernetes and Container Orchestration}
\label{sec:background-kubernetes}

Kubernetes~\cite{kubernetes} is the orchestration system used to deploy every
candidate in this thesis. It follows the desired-state model developed in the
Borg/Omega line of cluster-management systems~\cite{burns2016borg}: an
operator declares what should be running, and a control plane continuously
reconciles the cluster toward that state. The historical relationship is
useful context, but the thesis depends on Kubernetes' declarative and
resource-aware execution model rather than on Borg or Omega themselves.

Several Kubernetes primitives recur throughout the thesis. A \emph{Pod} is
the smallest deployable unit, containing one or more containers scheduled
together on a worker machine. A \emph{Deployment} manages a set of equivalent
Pod replicas and works to maintain the declared replica count. A \emph{Service}
provides a stable network identity and routes requests to healthy replicas.
A \emph{Namespace} groups resources into an isolated, disposable scope, which
the benchmark uses for each candidate deployment.

Pods also declare CPU and memory \emph{requests} and \emph{limits}. Requests
contribute to scheduler placement and capacity accounting; limits are enforced
runtime ceilings. A CPU limit can cause throttling, while exceeding a memory
limit can result in an out-of-memory termination. These distinctions matter
because a specification can be schedulable while still starving or throttling
the process under load.

The Kubernetes \emph{scheduler} turns a Deployment's replica count and
placement constraints into concrete worker-node assignments. It considers
each node's remaining allocatable capacity and any restrictions supplied by
the specification. Placement is therefore not merely an administrative
detail: it determines which Pods share a machine and which network paths they
use, and can change the performance observed by the load test.


% ----- §2.4  Deployment levers -----
\section{Deployment Configuration as a Performance Lever}
\label{sec:background-deployment-levers}

Once an application is deployed, its performance depends on both the code and
the configuration of the service around it. The agent in this thesis controls
a bounded set of deployment decisions; the exact fields and validation rules
are specified in Section~\ref{sec:methods-spec}. They can be understood as
four related families.

\paragraph{Application replicas.}
Replica count determines how many instances run behind the backend Service.
Too few replicas can leave requests queued behind busy instances. Additional
replicas can increase parallelism until another resource becomes the
bottleneck, but can also consume cluster capacity and multiply downstream
load, especially database connections.

\paragraph{Per-Pod resources.}
CPU and memory requests affect whether and where Pods can be scheduled;
limits constrain their runtime consumption. A request or limit that is too
small can cause starvation, throttling, or termination. Values that are much
larger than the workload requires waste allocatable capacity and may prevent
useful replicas or data-tier components from being scheduled. The relevant
quantity is therefore not a universally correct number, but the relationship
between the workload's demand and the cluster's finite capacity.

\paragraph{Placement.}
Placement specifies which workers a Pod may use and whether replicas should be
spread or allowed to co-locate. Co-location can make a small deployment
efficient, but it also makes replicas compete for the same CPU, memory, and
network path. Spreading replicas can reduce that contention while consuming
more of the cluster's placement budget. The best choice depends on the
workload and on what other Pods share the cluster.

\paragraph{Data-tier topology and connection handling.}
The backend can use a standalone database or a primary with read replicas. It
can connect directly or through a pooler that multiplexes application
connections onto fewer database connections; the deployment may also enable a
cache. These choices trade capacity, coordination overhead, and resource use.
A read replica cannot improve throughput if the application never sends reads
to it, and a pool or database connection ceiling can become the bottleneck
even when the application Pods have spare CPU.

These families are tunable without rewriting the application's source code,
which makes deployment configuration a distinct control surface. They are not
independent of code, however: connection-pool behavior, query patterns, and
application concurrency determine whether a deployment choice helps or hurts.
The benchmark therefore treats code and deployment as separate levers whose
effects must be observed in the same running system.


% ----- §2.5  Correctness vs performance -----
\section{Functional Correctness versus Runtime Performance}
\label{sec:background-correctness-vs-performance}

Functional correctness and runtime performance are related but different
properties. A functional test sends specified requests and checks their
responses against an expected contract: for example, whether an endpoint
returns the right object, enforces permissions, or rejects malformed input.
Passing such tests establishes correctness for the exercised behaviors. It
does not establish how quickly or reliably the service responds when many
requests overlap.

Conversely, a service can return responses quickly while computing the wrong
result. The load-test workload has its own request-level success and failure
marking, but that marking is not interchangeable with the separate BaxBench
functional-correctness verdict. This thesis therefore uses functional tests as
a gate before performance measurement, then uses concurrent load to measure
the running system.

Under load, a correct application can still fail because a query exhausts a
connection pool, a lock serialises requests, a database reaches its capacity,
or the application tier has too few replicas. Such failures often appear as
timeouts or elevated error rates, which is why performance feedback must be
interpreted together with application and cluster diagnostics.


% ----- §2.6  Measuring performance -----
\section{Measuring Performance Under Load}
\label{sec:background-goodput}

A \emph{load test} creates concurrent demand and records how a service responds
as that demand changes. In this thesis, Locust~\cite{locust} runs simulated
users that follow the scenario's verified request flows, so the load workload
is grounded in the same API contract used by the functional tests. The
workload defines what users request; the load profile controls how the number
of concurrent users changes over time.

Three quantities are central. \emph{Offered load} is the demand presented by
the clients. \emph{Throughput} is the rate of completed request attempts.
\emph{Latency} is the time from issuing a request to receiving its response,
usually summarised with percentiles such as p50 or p95. Percentiles matter
because a mean can hide a long tail of requests waiting behind a saturated
resource.

As offered load rises, some resource eventually becomes limiting: CPU, memory,
network capacity, database connections, or storage. The resulting
\emph{saturation point} is where additional offered load stops producing
proportionally more completed work. Near that point, queues grow and latency
often rises sharply; beyond it, requests may time out or fail rather than
merely become slower. The limiting resource is the deployment's bottleneck at
that load level.

Raw throughput is not sufficient as an objective because a system can complete
many requests while failing a substantial fraction of them. \textbf{Goodput}
is the rate of successful requests, excluding failed and timed-out requests.
In this thesis, a request is successful when the Locust workload marks it as
such; this is a request-level runtime signal, not a replacement for the
functional-correctness gate. The term is adopted from serving-system work that
optimises successful request rate under latency objectives
\cite{zhong2024distserve}, but here it is applied to ordinary HTTP backends.

The benchmark maximises goodput subject to health and stability conditions,
including failure-rate, latency, and settling checks. \emph{Sustained
goodput} is therefore an operational measurement: the best accepted stable
goodput level produced by the refinement phase of the load profile, rather
than a transient peak during the exploratory ramp. Section
\ref{sec:methods-load-profile} gives the exact phases, thresholds, and
aggregation procedure.


% ----- §2.7  Iterative refinement -----
\section{Iterative Refinement as an Agentic Improvement Pattern}
\label{sec:background-iterative-refinement}

Rather than treating generation as a single attempt, an iterative system
produces a candidate, evaluates it, feeds the outcome back to the model, and
produces a revised candidate. The model does not need to change its weights:
the history of previous candidates and observations becomes part of the next
input. This is a feedback loop over artifacts and measurements.

The stopping rule depends on the objective. A correctness-gated loop can stop
when a test suite passes. A performance-gated loop has no unique pass
condition, because a deployment can usually trade more iterations for a
possible improvement. It therefore stops when improvement stalls, a budget is
exhausted, or an external criterion is reached.

Feedback fidelity is equally important. Measurements taken in an environment
that resembles the target environment are less likely to reward a change that
works only in a proxy. This thesis consequently measures candidates on the
target Kubernetes cluster rather than using a local performance surrogate.
The general pattern is distinct from the particular prior systems that
instantiate it; those systems are reviewed in Chapter~\ref{ch:related-work}.


% ----- §2.8  Conceptual model -----
\section{Conceptual Model: Code and Deployment as Coupled Levers}
\label{sec:background-synthesis}

The thesis evaluates a candidate as a pair of artifacts: application code
$c$ and deployment configuration $d$. For a fixed workload $w$ and cluster
capacity $C$, the running system has an observed performance outcome that can
be written conceptually as
\[
  G(c,d;w,C),
\]
where $G$ is sustained goodput under the benchmark's measurement procedure.
The candidate must also satisfy feasibility conditions: the code must pass
the functional gate, the deployment must be schedulable and ready, and the
load test must meet its health and stability conditions.

Code and deployment are distinct control surfaces. Changing source code can
alter query patterns, concurrency, or connection usage without changing the
replica count; changing the specification can change replicas, resources, or
topology without changing what the application computes. They are nevertheless
coupled through the running system. A code defect can look like a deployment
shortfall, and a capacity bottleneck can look like an application failure.
Consequently, the scalar goodput score is interpreted together with latency,
failure, resource, and database diagnostics.

There is also no deployment configuration that is optimal independently of its
context. The useful choice of $d$ depends on the workload mix, the application
framework, and the available cluster capacity $C$. The evaluation keeps those
factors fixed within an experiment task and asks whether feedback allows the
agent to improve the particular system it is running. Chapter~\ref{ch:methods}
now specifies how this conceptual model is operationalised through artifact
validation, deployment, load measurement, and one-lever-at-a-time refinement.
